\documentclass[11pt]{article}

\usepackage[T1]{fontenc}
\usepackage[utf8]{inputenc}
\usepackage{lmodern}
\usepackage[a4paper,margin=1in]{geometry}
\usepackage{amsmath}
\usepackage{amssymb}
\usepackage{booktabs}
\usepackage{xcolor}
\usepackage{enumitem}
\usepackage{fancyhdr}
\usepackage[hidelinks]{hyperref}
\usepackage{titlesec}

% --- cores ---
\definecolor{ink}{HTML}{1A1A2E}
\definecolor{klein}{HTML}{1F4FD6}
\definecolor{soft}{HTML}{F4F5F7}
\definecolor{rule}{HTML}{C9CED6}
\definecolor{gold}{HTML}{B8860B}

\hypersetup{colorlinks=true, linkcolor=klein, urlcolor=klein, citecolor=klein}

% --- títulos ---
\titleformat{\section}{\Large\bfseries\color{ink}}{\thesection}{0.6em}{}
\titleformat{\subsection}{\large\bfseries\color{ink}}{\thesubsection}{0.6em}{}
\titlespacing*{\section}{0pt}{1.4em}{0.6em}

% --- cabeçalho / rodapé ---
\pagestyle{fancy}
\fancyhf{}
\renewcommand{\headrulewidth}{0.4pt}
\fancyhead[L]{\small\color{ink}\textbf{SPP Compliance Layer}}
\fancyhead[R]{\small\color{ink}paper de negócio}
\fancyfoot[C]{\small\thepage}

% --- caixa de destaque ---
\newcommand{\keybox}[1]{%
  \par\vspace{0.6em}%
  \noindent\colorbox{soft}{\parbox{\dimexpr\linewidth-2\fboxsep}{%
  \vspace{0.3em}\small #1\vspace{0.3em}}}%
  \par\vspace{0.6em}%
}

\newcommand{\code}[1]{\texttt{\small #1}}

\setlist[itemize]{leftmargin=1.3em, itemsep=0.25em, topsep=0.35em}

\begin{document}

% --- capa ---
\begin{center}
{\LARGE\bfseries\color{ink} SPP Compliance Layer}\\[0.4em]
{\large\color{klein} Capturar capital institucional para a Stellar}\\[0.2em]
{\normalsize\color{ink} fechando o trilema privacidade $\times$ compliance $\times$ durabilidade quântica}\\[1.2em]
{\small Manuel Galmanus \quad$\cdot$\quad 07 de agosto de 2026}
\end{center}

\vspace{0.6em}
\hrule height 0.4pt
\vspace{1.2em}

\section*{Sumário executivo}
\addcontentsline{toc}{section}{Sumário executivo}

Capital institucional não entra em blockchain pública por três razões que se
bloqueiam mutuamente. Uma tesouraria não pode expor posições e contrapartes
(\textbf{privacidade}); um regulador não aceita privacidade sem allow-list e
trilha de auditoria (\textbf{compliance}); e nenhum dos dois aceita que a prova de
compliance se apoie em criptografia que um computador quântico vai quebrar num
livro-razão permanente (\textbf{durabilidade}). A Stellar já resolveu a primeira
perna, com os primitivos de pagamento privado da Nethermind e o Confidential Token
da OpenZeppelin. As duas outras estão abertas --- e são exatamente onde o
\textbf{SPP Compliance Layer} atua.

\keybox{\textbf{Tese.} O SPP Compliance Layer transforma a privacidade que a
Stellar já tem no produto que uma instituição pode de fato usar: uma política de
compliance provada honesta por criptografia \textbf{pós-quântica} e imposta
\textbf{on-chain}. Com isso, remove o bloqueador que hoje mantém tesourarias,
custodiantes e emissores de RWA fora da rede.}

O ponto não é ``mais uma prova de conhecimento zero''. Os primitivos de
privacidade da Stellar hoje provam com \emph{pairing} sobre BN254 (Groth16,
UltraHonk), e a própria SDF afirma, no Quantum Preparedness Plan:
\emph{``there is no drop-in post-quantum replacement for pairing-based SNARKs.''}
Uma atestação de compliance construída sobre essa curva pode, com o tempo, ser
\textbf{forjada}. O SPP Compliance Layer atesta a parte que se pode tornar
pós-quântica hoje --- o histórico da política de compliance --- com um STARK
baseado em hash, verificado dentro de um contrato Soroban.

\section{O problema, com precisão}

Um pagamento privado numa \emph{privacy pool} esconde valor e contraparte. Para
uma instituição isso é condição necessária, não suficiente. Faltam três coisas, e
é a combinação delas que trava a migração:

\begin{enumerate}[leftmargin=1.6em, itemsep=0.4em]
\item \textbf{A política de compliance é esquecida pela rede.} A carteira
reconstrói estado relendo todos os eventos que a \emph{pool} já emitiu; o Soroban
RPC guarda apenas $\approx$7 dias desses eventos (janela medida: 120.959 ledgers
$\approx$ 7,01 dias, que ainda \emph{desliza} com a cadeia). O histórico das
raízes de compliance --- o \emph{association set} que um provedor (ASP) publica ---
some da RPC durante um fim de semana. Sem esse histórico, ninguém prova, depois,
que a política era íntegra.

\item \textbf{A lacuna já tem um remédio oficial --- inseguro.} O cliente SPP da
Nethermind já faz \emph{handoff} para um \code{bootnode\_url} quando bate no buraco
de retenção (\code{sdk/client/src/sync.rs}). E a documentação da própria Nethermind
(\code{docs/src/bootnode.md}) \textbf{nomeia} \emph{forged-history},
\emph{selective-omission} e \emph{misleading-handoff} como riscos \emph{não}
mitigados --- a única mitigação oferecida é ``cruze com múltiplos provedores de
RPC'', que é confiança social, não prova.

\item \textbf{A prova de honestidade é quebrável por Shor.} Mesmo com o histórico
preservado, atestá-lo com Groth16/BN254 significa um \emph{trusted setup} para
confiar e um adversário quântico para sobreviver, num registro permanente.
``Harvest now, forge later'' vale para atestações tanto quanto para cifragem.
\end{enumerate}

Uma tesouraria olha para isso e conclui, corretamente, que a privacidade da
Stellar ainda não é utilizável em escala institucional. Não por falta de
privacidade --- por falta da camada que a torna auditável, durável e imposta.

\section{A solução, concretamente}

Três peças fecham exatamente essas três lacunas, e o laço fecha nos dois lados:
sem prova, nenhuma política admitida; sem política admitida, nenhum gasto.

\begin{center}
\small
\begin{tabular}{@{}p{3.1cm}p{6.2cm}p{4.4cm}@{}}
\toprule
\textbf{Peça} & \textbf{O que impõe} & \textbf{Custo on-chain medido} \\
\midrule
Atestação pós-quântica & prova que o histórico de raízes de compliance é uma
sequência \emph{append-only} sem lacunas (Circle-STARK baseado em hash, sem
\emph{trusted setup}, sem \emph{pairing}) & \emph{proof} $\approx$131 KB
(off-chain, 62 bits quânticos) \\
\addlinespace
Gate de compliance \code{admit\_root} & registra uma raiz como válida on-chain
\textbf{só se} a prova pós-quântica verificar na mesma transação & 260M de 400M
instruções (65\% do teto), wasm 115 KB, $\approx$0,0385 XLM \\
\addlinespace
Pool com política \code{spend} & liquida um pagamento confidencial \textbf{só
contra uma raiz admitida} (\emph{cross-contract} \code{is\_attested}) & recibo
on-chain, \emph{replay} recusado \\
\bottomrule
\end{tabular}
\end{center}

Recibos vivos na testnet (clicáveis no repositório): a raiz é admitida na
transação \code{a2c3227c\ldots}, o gasto liquida contra ela na
\code{255db58d\ldots}, e uma prova adulterada ou raiz não-admitida é
\textbf{recusada on-chain}. Nada é mockado: cada número sai de um comando do
próprio repositório.

O diferencial técnico contra o estado da arte é estreito e verificável: todo
verificador de conhecimento zero on-chain na Stellar hoje é \emph{pairing-based}
(Groth16, UltraHonk, RISC~Zero via seu embrulho Groth16). Este é --- até onde a
busca pública alcança --- o primeiro STARK \textbf{transparente, baseado em hash},
verificado \textbf{nativamente} num contrato Soroban e tornado \emph{load-bearing}
numa política de compliance. Contra um adversário quântico ele apenas
\emph{enfraquece} (Grover corta os bits pela metade); nunca \emph{quebra} (Shor).
É a diferença entre ``enfraquecido, some mais consultas'' e ``quebrado, comece de
novo''.

\section{Por que isto destrava migração de capital}

O argumento causal, sem retórica:

\begin{itemize}
\item Uma instituição só move capital para um trilho quando os três ---
confidencialidade, compliance e durabilidade --- estão satisfeitos \textbf{ao
mesmo tempo}. Falha em um veta o conjunto.
\item A Stellar entrega o primeiro com os primitivos existentes. O SPP Compliance
Layer entrega os outros dois: índice durável $+$ \emph{bootnode} verificável
(fecha o buraco de retenção e os riscos que a Nethermind documenta como abertos) e
atestação pós-quântica imposta on-chain (fecha o risco de forja futura).
\item Portanto o Compliance Layer não compete com as carteiras da rede --- \textbf{é
aquilo de que elas dependem}. Cada carteira aponta seu \code{bootnode\_url} para cá
e passa a \emph{gatear} compliance na mesma atestação pós-quântica.
\end{itemize}

\keybox{\textbf{Posicionamento, com disciplina.} Isto não é uma custódia nem um
trilho de liquidação com garantia de recuperação de fundos. É a camada de
\emph{prova de integridade da política} sobre a qual uma custódia pode se apoiar.
Torna a privacidade da Stellar auditável e durável; não guarda o ativo nem assume
obrigação de VASP. Esse enquadramento é o que sustenta a conversa com um jurídico
institucional.}

\section{Por que Stellar, e não Ethereum ou Solana}

Um argumento de plataforma, com número, não com torcida:

\begin{itemize}
\item O mesmo relacionamento criptográfico já rodou nas três redes. Em Solana a
parede foi o \emph{heap} de 256 KB (2,38M CU em LiteSVM). Na Stellar a memória
pica em 2,1 MB de 40 MB e a parede de CPU caiu para 62\% do teto de uma
transação, \textbf{sem mudança de protocolo} --- só uma \emph{host-function} de
keccak que já existe e uma \emph{flag} de compilador.
\item O que no Ethereum L1 exigiu um verificador Solidity dedicado e anos de
engenharia de uma empresa feita para isso, aqui roda como um contrato wasm comum
de $\approx$115 KB que qualquer um implanta, em uma transação, por $\approx$2
centavos de dólar.
\item Ou seja: a Stellar tem, hoje e sem \emph{fork}, o \emph{headroom} para
carregar criptografia pós-quântica pesada dentro dos limites de uma transação.
Esse é um argumento concreto para trazer a migração \emph{para cá}.
\end{itemize}

\section{Mercado: a dor tem dono, e tem número}

O valor de RWA tokenizado on-chain (ex-stablecoins) saiu de $\approx$US\$~6~bi no
início de 2025 para \textbf{US\$~33,5~bi} em julho de 2026 (rwa.xyz, via Stobox),
um crescimento de $\approx$4--5$\times$ em dezoito meses; tesouros tokenizados são
$\approx$2/3 do setor. As projeções para 2030 vão de US\$~2~tri (McKinsey, base
conservadora) a US\$~16~tri (BCG) --- um intervalo de 8$\times$ entre analistas
sérios, em que a direção é consenso e só a inclinação é disputada.

A Stellar não está de fora dessa corrida: é a \textbf{quarta maior rede de RWA,
com 9,0\%} do valor on-chain (Stobox, jul/2026), atrás de Ethereum, BNB e Solana;
foi na Stellar que a Franklin Templeton lançou seu fundo on-chain (BENJI). O
capital já está chegando --- e a próxima leva depende de infraestrutura, não de
mais permissão regulatória.

Porque o gargalo declarado, por quem move o dinheiro, é privacidade e compliance:

\begin{center}
\small
\begin{tabular}{@{}p{3.6cm}p{9.6cm}@{}}
\toprule
\textbf{Fonte} & \textbf{O que diz} \\
\midrule
Visa, jun/2026 & ``Many banks see the lack of privacy as a dealbreaker for moving
meaningful activity onchain.'' (Rubail Birwadker, Global Head of Growth Products) \\
\addlinespace
DRW, mar/2026 & Publicar toda operação on-chain seria ``a failure of fiduciary
duty''; ``privacy is kind of at the top of the list'' para adoção institucional.
(Don Wilson, CEO) \\
\addlinespace
EEA Privacy WG & ``Privacy and confidentiality are not features --- they are
prerequisites for enterprise adoption.'' \\
\addlinespace
Fireblocks, abr/2026 & O que falta é ``the compliance layer alongside the privacy
layer, which is where most implementations fall short.'' \\
\bottomrule
\end{tabular}
\end{center}

E o diagnóstico de infraestrutura é explícito: relatórios de 2026 (BitMart,
RedStone) apontam que \emph{``the deployment gap is an infrastructure problem, not
a regulatory one''} --- de $\approx$US\$~27~bi em RWA tokenizado, apenas
$\approx$10\% é ativamente utilizável; os outros $\approx$90\% ficam parados em
carteiras, rendendo mas sem circular.

\keybox{A leitura de mercado converge para uma frase: falta a \textbf{camada de
compliance ao lado da camada de privacidade}. É exatamente o que este projeto é
--- com uma diferença que ninguém mais no setor está entregando: a atestação de
compliance é \textbf{pós-quântica e imposta on-chain}, não pairing-based.}

\subsection*{Contexto competitivo, sem ilusão}

Parte do capital institucional escolhe redes permissionadas (Canton, com a Visa
como \emph{Super Validator}; o Kinexys do JPMorgan liquida US\$~7~bi/dia). O SPP
Compliance Layer não disputa esse segmento fechado. Ele atende o capital que
escolhe trilhos \textbf{confidenciais em cadeia pública} --- o campo que Visa,
Chainlink (Privacy Standard, prova de compliance on-chain) e Fireblocks estão
correndo para equipar, e onde a Stellar já tem 9\% do RWA. A aposta é estreita e
explícita: nesse campo, ser pós-quântico e nativo do Soroban é uma vantagem que os
concorrentes pairing-based não têm.

\section{Público-alvo}

\begin{itemize}
\item \textbf{Emissores de RWA e ativos tokenizados regulados} --- precisam de
allow-list imposta e de trilha de auditoria que sobreviva à vida longa do ativo.
\item \textbf{Tesourarias e mesas de OTC/B2B} --- pagamento confidencial entre
partes com política de identidade, sem expor posição ao mercado.
\item \textbf{Custodiantes e provedores de compliance (ASPs)} --- operam a
política; o Compliance Layer dá a eles uma atestação que um regulador aceita como
prova, não como promessa.
\item \textbf{Reguladores e auditores} --- o consumidor final da atestação:
verificam a integridade de um conjunto de aprovação de 2026 em 2035, mesmo depois
que um quântico poderia forjar um BN254.
\end{itemize}

O gancho comum a todos: precisam de privacidade \emph{do mercado}, não do
regulador --- confidencialidade seletiva com trilha de auditoria --- e precisam
que essa trilha sobreviva à vida longa do ativo. É o par exato que a seção de
mercado documenta como o gargalo aberto.

\section{Como monetiza}

Três linhas de receita, em ordem de proximidade:

\begin{enumerate}[leftmargin=1.6em]
\item \textbf{Bootnode verificável como serviço.} Operar o \emph{bootnode}
confiável que os clientes SPP apontam, com SLA e prova de cobertura. Assinatura
por instituição.
\item \textbf{Atestação como serviço.} Cobrar por atestação pós-quântica de
histórico de compliance (por período/por \emph{pool}) para ASPs, emissores e
custodiantes que precisam da prova para relatório regulatório.
\item \textbf{Licenciamento do verificador on-chain $+$ SDK.} O contrato-\emph{gate}
e o SDK de integração para quem quer \emph{gatear} liquidação na própria política.
\end{enumerate}

O primeiro marco de negócio é uma integração real --- um piloto com um ASP,
emissor ou custodiante --- não um preço de tabela.

\section{Roadmap}

\begin{itemize}
\item \textbf{Agora:} testnet, laço completo provado (\emph{gate} $+$ \emph{pool}),
recibos vivos, dupla submissão no Stellar Summit (raia \emph{Wallets} e raia
\emph{Enterprise/Compliance/RWA}).
\item \textbf{Curto prazo:} elevar a segurança quântica on-chain de 24 bits (no
orçamento de uma transação, 40 consultas) via camada recursiva --- a técnica
STARKPack do próprio patrocinador --- a custo limitado. Off-chain já são 62 bits.
\item \textbf{Médio prazo:} \emph{deploy} em mainnet (os limites de mainnet foram
lidos de \code{getLedgerEntries} e todos os números cabem identicamente --- é
decisão de \emph{funding}, não de engenharia). Primeira integração com um
ASP/emissor real.
\item \textbf{Longo prazo:} \emph{wiring} da \emph{pool} para \emph{exigir} a
atestação como pré-condição de liquidação.
\end{itemize}

\section{Considerações de implementação}

Ditas na frente, porque é o que sustenta credibilidade técnica --- cada uma com
seu endereçamento:

\begin{itemize}
\item \textbf{A atestação prova estrutura de índice, não legitimidade de raiz.}
Ela prova que o histórico é uma sequência \emph{append-only} sem lacunas sobre
raízes testemunhadas, com \emph{endpoints} fixados em valores públicos --- legíveis
on-chain. A ligação de \emph{completude} é a prova de cobertura da camada de índice.
As duas juntas dão a garantia; declaradas com precisão, não infladas.
\item \textbf{24 bits quânticos on-chain} no orçamento de uma transação, contra 62
off-chain. O caminho recursivo (STARKPack) é o plano de escala para elevar a cifra
on-chain a custo limitado.
\item \textbf{Testnet, com caminho de mainnet medido.} Os limites de mainnet já
foram lidos e todos os números cabem; o \emph{deploy} é decisão de \emph{funding}.
\item \textbf{Enquadramento jurídico.} O posicionamento correto --- camada de prova
de integridade, não custódia --- mantém o produto fora da obrigação de VASP e é
sustentado com disciplina em todo material externo.
\end{itemize}

\section*{Fecho}
\addcontentsline{toc}{section}{Fecho}

A Stellar já tem privacidade. O que falta para ela receber capital institucional
não é mais privacidade --- é a camada que torna essa privacidade auditável por um
regulador e durável através da transição quântica, imposta pelo próprio
livro-razão. Essa camada existe, roda na testnet com recibos, e é honesta sobre
onde para. O trabalho de negócio a partir daqui é uma coisa só: converter uma
prova técnica que funciona em uma integração institucional que paga.

\vspace{1em}
\hrule height 0.4pt
\vspace{0.6em}
\noindent{\small\color{ink} Repositório: \href{https://github.com/Galmanus/spp-compliance-layer}{github.com/Galmanus/spp-compliance-layer} \quad$\cdot$\quad MIT \quad$\cdot$\quad recibos e verificação independente em \code{VERIFY-IT-YOURSELF.md}}

\section*{Fontes de mercado}
\addcontentsline{toc}{section}{Fontes de mercado}
\footnotesize
\begin{itemize}[leftmargin=1.3em, itemsep=0.15em]
\item Stobox, \emph{State of RWA Tokenization --- 2026 Mid-Year Report}, jul/2026 (US\$~33,5~bi on-chain; Stellar 9,0\%). \href{https://www.stobox.io/reports/state-of-rwa-2026}{stobox.io/reports/state-of-rwa-2026}
\item CoinGecko, \emph{RWA Report 2026} (US\$~5,42~bi $\to$ US\$~19,3~bi, $+$256,7\%). \href{https://www.coingecko.com/research/publications/rwa-report-2026}{coingecko.com}
\item BitMart / RedStone, \emph{State of RWA}, mai/2026 (US\$~24,6~bi TVL; ``infrastructure problem, not regulatory''; $\approx$10\% composável). \href{https://www.bitmart.com/static-file/resources/public/BitMart_State_of_RWA_Issue02_May2026.pdf}{bitmart.com}
\item Chainalysis, \emph{Tokenized RWAs}, abr/2026 ($\approx$US\$~30~bi AUM; capital institucional supera o de varejo).
\item Visa / FinTech Magazine, jun/2026 (``lack of privacy as a dealbreaker''). \href{https://fintechmagazine.com/news/visa-anchoring-trust-in-next-generation-enterprise-payments}{fintechmagazine.com}
\item CoinDesk, \emph{DRW's Don Wilson on open ledgers}, mar/2026 (``failure of fiduciary duty''). \href{https://www.coindesk.com/business/2026/03/26/why-big-banks-are-snubbing-open-ledgers-to-build-their-own-private-blockchains}{coindesk.com}
\item Fireblocks (Ran Goldi), \emph{The Blockchain Privacy Problem}, abr/2026 (``compliance layer alongside the privacy layer''). \href{https://www.fireblocks.com/blog/blockchain-privacy-problem}{fireblocks.com}
\item Enterprise Ethereum Alliance, \emph{State of Privacy on Ethereum for Enterprise} (``prerequisites for enterprise adoption'').
\item Projeções 2030: McKinsey (US\$~2~tri, base) e BCG (US\$~16~tri).
\end{itemize}

\end{document}
